\documentclass[a4paper,12pt]{article}
\usepackage[T2A]{fontenc}
\usepackage[utf8]{inputenc}
\usepackage[russian]{babel}
\usepackage{cmap}
\usepackage{amsmath,amssymb,mathtools}
\usepackage{geometry}
\geometry{left=24mm,right=24mm,top=23mm,bottom=24mm}
\usepackage{microtype}
\usepackage{icomma}
\usepackage{tikz}
\usetikzlibrary{arrows.meta,positioning,shapes.geometric,calc}
\usepackage{tabularx,booktabs}
\usepackage{enumitem}
\usepackage{parskip}
\usepackage{fancyhdr}
\usepackage{setspace}
\usepackage{xcolor}
\usepackage{hyperref}
\definecolor{accent}{HTML}{007B80}
\definecolor{darkaccent}{HTML}{00535D}
\definecolor{textgray}{HTML}{34383B}
\definecolor{softgray}{HTML}{778186}
\hypersetup{colorlinks=true,linkcolor=darkaccent}
\setlength{\parindent}{0pt}
\setlength{\parskip}{4.5pt}
\setstretch{1.17}
\setlist[enumerate]{topsep=5pt,itemsep=5pt,leftmargin=*,label=\textcolor{accent}{\arabic*.}}
\setlist[itemize]{topsep=4pt,itemsep=4pt,leftmargin=*}
\newcommand{\sectionrule}{\par\vspace{3pt}{\color{accent}\hrule height .45pt}\vspace{9pt}}
\newcommand{\subidea}[1]{\vspace{5pt}{\color{darkaccent}\large\bfseries #1}\par\vspace{2pt}}
\newcommand{\kmh}{\text{ км/ч}}
\newcommand{\km}{\text{ км}}
\pagestyle{fancy}
\fancyhf{}
\fancyhead[L]{\footnotesize\textcolor{softgray}{Задачи на движение}}
\fancyhead[R]{\footnotesize\textcolor{softgray}{Ярослав Верещагин}}
\fancyfoot[C]{\footnotesize\textcolor{softgray}{\thepage}}
\renewcommand{\headrulewidth}{.4pt}
\renewcommand{\headrule}{\hbox to\headwidth{\color{accent}\leaders\hrule height \headrulewidth\hfill}}
\setlength{\headheight}{15pt}
\newcommand{\sect}[1]{\section{#1}\sectionrule}
\renewcommand{\thesection}{\arabic{section}}
\begin{document}
\thispagestyle{empty}
\begin{center}
\vspace*{27mm}
{\color{darkaccent}\Large МАТЕМАТИКА}\par
\vspace{24mm}
{\fontsize{25}{30}\selectfont Чек-лист}\par
\vspace{14mm}
{\color{accent}\fontsize{22}{27}\selectfont Задачи на движение:\par
встреча и догонка}\par
\vspace{23mm}
{\large Подготовка к ОГЭ по математике}\par
\vfill
{\large Ярослав Верещагин}\par
\vspace{9mm}
{\large 2 ноября 2026 г.}\par
\vspace{16mm}
{\small Лёвин Артём Александрович\par
эксклюзивно для Ярослава Верещагина}\par
\vspace{21mm}
\begin{tikzpicture}
\draw[accent,line width=.9pt] (0,0) .. controls (1.5,.6) and (2.9,-.6) .. (4.2,0)
 .. controls (5.8,.7) and (7.1,-.6) .. (8.5,0);
\fill[accent] (4.2,0) circle (1.4pt);
\end{tikzpicture}
\end{center}
\clearpage

\sect{Главная формула и таблица}
В задачах на движение связываем три величины: скорость $v$, время $t$ и пройденный путь $S$.
\[
\boxed{S=vt}\qquad \boxed{v=\frac{S}{t}}\qquad \boxed{t=\frac{S}{v}}.
\]
Если скорость выражена в километрах в час, а время --- в часах, путь получится в километрах. При делении на время предполагается $t>0$, при делении на скорость --- $v>0$.

\subidea{Как заполнять таблицу}
\begin{enumerate}
\item \textbf{Выберите неизвестное $x$.} Удобнее обозначить через $x$ именно то, что просят найти.
\item \textbf{Запишите данные условия.} Для каждого участника укажите известные скорость, время и путь.
\item \textbf{Вычислите оставшиеся клетки по формуле.} Сначала заполните таблицу, затем используйте информацию о встрече или разнице во времени.
\item \textbf{Составьте уравнение.} Сравнивайте величины одной размерности: километры с километрами, часы с часами.
\end{enumerate}

\sect{Встречное движение: общий момент выезда}
Два автомобиля одновременно выезжают из городов $A$ и $B$ навстречу друг другу. До встречи они едут \textbf{одинаковое время}; вместе проходят всё расстояние между городами:
\[
 t_1=t_2,\qquad S_1+S_2=D.
\]
\subidea{Пример 1. Нужно найти время}
Расстояние между городами $560$ км; скорости $60$ и $80$ км/ч. Пусть $x$ --- число часов до встречи.
\begin{center}
\begin{tabular}{@{}lccc@{}}\toprule
Участник & $v$, км/ч & $t$, ч & $S$, км\\\midrule
Из $A$ & $60$ & $x$ & $60x$\\
Из $B$ & $80$ & $x$ & $80x$\\\bottomrule
\end{tabular}
\end{center}
\[
60x+80x=560,\quad 140x=560,\quad \boxed{x=4\text{ ч}}.
\]
Главная подсказка: при одновременном выезде неизвестное \emph{время} можно разместить в обеих строках.

\clearpage
\sect{Когда просят найти расстояние или скорость}
\subidea{Пример 2. Нужно найти путь одного автомобиля}
Из городов, расстояние между которыми $405$ км, одновременно выехали навстречу автомобили со скоростями $60$ и $75$ км/ч. Найдите расстояние от города $B$ до места встречи.

Пусть $x$ км --- путь автомобиля из $B$. Тогда путь автомобиля из $A$ равен $405-x$.
\begin{center}
\begin{tabular}{@{}lccc@{}}\toprule
Участник & $v$, км/ч & $t$, ч & $S$, км\\\midrule
Из $A$ & $60$ & $\dfrac{405-x}{60}$ & $405-x$\\[5pt]
Из $B$ & $75$ & $\dfrac{x}{75}$ & $x$\\[5pt]\bottomrule
\end{tabular}
\end{center}
Время одинаковое:
\[
\frac{405-x}{60}=\frac{x}{75}\quad\Longrightarrow\quad
75(405-x)=60x\quad\Longrightarrow\quad \boxed{x=225\km}.
\]
Проверка: $180/60=225/75=3$ часа. Ограничение: $0<x<405$.

\subidea{Пример 3. Нужно найти скорость}
Два автомобиля одновременно двигались навстречу. Между городами $540$ км, скорость первого $70$ км/ч, встреча произошла через $4$ часа. Скорость второго равна $x$ км/ч.
\[
\underbrace{70\cdot4}_{\text{первый путь}}+
\underbrace{4x}_{\text{второй путь}}=540,
\qquad 4x=260,\qquad \boxed{x=65\kmh}.
\]
\subidea{Два равноправных способа записать уравнение}
При встречном движении можно использовать разные данные:
\begin{itemize}
\item \textbf{Сумму путей:} $S_1+S_2=D$ (если известен весь путь между городами).
\item \textbf{Равенство времён:} $\dfrac{S_1}{v_1}=\dfrac{S_2}{v_2}$ (если выезд одновременный).
\end{itemize}
Не нужно искусственно использовать обе связи в одном уравнении. Выбирайте ту, которая позволяет выразить неизвестное проще.

\clearpage
\sect{Встречное движение: разное время выезда}
Если второй автомобиль выехал через $\Delta t$ часов после первого, то первый к моменту встречи находился в пути дольше:
\[
\boxed{t_1-t_2=\Delta t},\qquad t_1>t_2>0.
\]
\subidea{Пример 4. Искомое --- время первого}
Расстояние между городами $645$ км. Первый автомобиль едет со скоростью $70$ км/ч. Через $3$ часа навстречу ему выехал второй со скоростью $75$ км/ч. Через сколько часов \emph{после выезда первого} они встретятся?
\begin{center}
\begin{tabular}{@{}lccc@{}}\toprule
Участник & $v$, км/ч & $t$, ч & $S$, км\\\midrule
Первый & $70$ & $x$ & $70x$\\
Второй & $75$ & $x-3$ & $75(x-3)$\\\bottomrule
\end{tabular}
\end{center}
\[
70x+75(x-3)=645\quad\Longrightarrow\quad\boxed{x=6\text{ ч}}.
\]
Проверка: $70\cdot6+75\cdot3=420+225=645$ км. Обязательно $x>3$.

\subidea{Пример 5. Искомое --- время второго}
Расстояние $580$ км. Первый автомобиль движется со скоростью $80$ км/ч, второй --- $60$ км/ч и выезжает на $2$ часа позже. Сколько часов ехал второй?
\[
 t_2=x,\qquad t_1=x+2,
\]
\[
80(x+2)+60x=580\quad\Longrightarrow\quad\boxed{x=3\text{ ч}}.
\]
\subidea{Пример 6. Когда известны пути, сравните времена}
Города разделяет $380$ км. Автомобиль из $A$ едет со скоростью $50$ км/ч, из $B$ --- $60$ км/ч. Второй выехал на час позже. Пусть $x$ км --- расстояние от $A$ до встречи.
\[
\underbrace{\frac{x}{50}}_{t_1}-
\underbrace{\frac{380-x}{60}}_{t_2}=1,
\qquad \boxed{x=200\km}.
\]
Это \emph{разность времени}, поэтому справа стоит $1$ час, а не $380$ км.

\clearpage
\sect{Задача на догонку}
Оба участника движутся \textbf{в одну сторону}. Более быстрый стартует на $d$ километров позади другого. При одновременном начале движения их времена одинаковы, а быстрый проходит \emph{лишние} $d$ километров:
\[
\boxed{S_{\text{быстрый}}=d+S_{\text{медленный}}}.
\]
\begin{center}
\begin{tikzpicture}[x=1cm,y=1cm,>=Stealth]
\draw[line width=1pt,textgray] (0,0)--(11.3,0);
\draw[line width=1.3pt,accent,{Stealth}-{Stealth}] (0,-.4)--node[below]{отставание $d=125$ км}(3.3,-.4);
\draw[line width=1.1pt,darkaccent,-{Stealth[length=3mm]}] (.15,.4)--(10.6,.4);
\draw[line width=1.1pt,accent,-{Stealth[length=3mm]}] (3.3,.72)--(10.6,.72);
\node[below] at (0,0) {$A$};
\node[below] at (3.3,0) {$B$};
\node[below] at (10.6,0) {$C$};
\node[above,anchor=south west,darkaccent] at (4,.4) {легковой автомобиль};
\node[above,anchor=south west,accent] at (5.8,.72) {грузовик};
\end{tikzpicture}
\end{center}
\subidea{Пример 7. Скорость больше на 25 км/ч}
Легковой автомобиль выехал из $A$, грузовик --- из $B$, расположенного на $125$ км впереди. Стартовали одновременно, двигались в одном направлении. Легковой автомобиль на $25$ км/ч быстрее. Через сколько часов он догонит грузовик?

Обозначим скорость грузовика через $v$ км/ч, время до встречи через $x$ часов.
\begin{center}
\begin{tabular}{@{}lccc@{}}\toprule
Участник & $v$, км/ч & $t$, ч & $S$, км\\\midrule
Легковой & $v+25$ & $x$ & $(v+25)x$\\
Грузовик & $v$ & $x$ & $vx$\\\bottomrule
\end{tabular}
\end{center}
\[
(v+25)x=125+vx\quad\Longrightarrow\quad25x=125
\quad\Longrightarrow\quad\boxed{x=5\text{ ч}}.
\]
Время догонки зависит от \textbf{разности скоростей}: $x=d/(v_{\text{быстрый}}-v_{\text{медленный}})$. Здесь знаменатель положителен.

\subidea{Типичные ошибки}
\begin{itemize}
\item При позднем выезде прибавляют задержку к \emph{пути}; задержка измеряется часами и должна участвовать во времени.
\item Складывают часы и приравнивают результат к километрам. Проверяйте единицы обеих частей уравнения.
\item Если $x$ --- время второго, ошибочно пишут $x-\Delta t$ для первого. Раньше выехавший движется дольше, значит, его время $x+\Delta t$.
\item В догонке складывают пройденные пути. Правильная связь учитывает начальный отрыв.
\end{itemize}

\clearpage
\begin{center}
{\color{darkaccent}\LARGE Алгоритм: от текста к уравнению}\par\vspace{5mm}
\end{center}
\begin{center}
\begin{tikzpicture}[>=Stealth,
 node distance=11mm,
 start/.style={ellipse,draw=darkaccent,very thick,align=center,minimum width=44mm,minimum height=9mm,fill=white},
 processnode/.style={rounded corners=2mm,draw=accent,line width=.8pt,align=center,text width=104mm,inner sep=6pt},
 decision/.style={diamond,draw=darkaccent,line width=.9pt,aspect=2.6,align=center,text width=41mm,inner sep=2pt},
 branch/.style={rounded corners=2mm,draw=accent,line width=.8pt,align=center,text width=48mm,minimum height=25mm,inner sep=5pt},
 arr/.style={-{Stealth[length=3mm]},line width=.85pt,textgray}]
\node[start] (begin) {Прочитайте условие};
\node[processnode,below=of begin] (var) {Обозначьте через $x$ искомую величину\\\small (время, путь или скорость)};
\node[processnode,below=of var] (table) {Составьте таблицу $v$--$t$--$S$.\\Заполните оставшиеся клетки по формуле $S=vt$};
\node[decision,below=16mm of table] (which) {Как движутся\\ участники?};
\node[branch,below left=17mm and 3mm of which] (toward) {\textbf{Навстречу}\\[3pt]$S_1+S_2=D$\\[3pt]Времена равны или отличаются на задержку};
\node[branch,below right=17mm and 3mm of which] (catch) {\textbf{Вдогонку}\\[3pt]$S_{\text{быстр.}}=d+S_{\text{медл.}}$\\[3pt]Учитывайте начальный отрыв};
\node[processnode,below=61mm of which] (finish) {Запишите уравнение. Проверьте единицы измерения, ограничения и ответ};
\node[start,below=of finish] (end) {Готово};
\draw[arr] (begin)--(var);
\draw[arr] (var)--(table);
\draw[arr] (table)--(which);
\draw[arr] (which.south west)--node[above left,font=\small] {навстречу} (toward.north);
\draw[arr] (which.south east)--node[above right,font=\small] {вдогонку} (catch.north);
\draw[arr] (toward.south) |- (finish.west);
\draw[arr] (catch.south) |- (finish.east);
\draw[arr] (finish)--(end);
\end{tikzpicture}
\end{center}

\clearpage
\section*{Тренировочный блок --- Путь А}
\sectionrule
Для каждого задания \textbf{составьте таблицу «скорость -- время -- путь» и уравнение}. Затем найдите требуемую величину. Располагайте данные в таблице согласно условию, а не в произвольном порядке.
\begin{enumerate}
\item Два автомобиля одновременно выехали из двух городов навстречу друг другу. Расстояние между городами $480$ км, скорости $55$ и $65$ км/ч. Через сколько часов они встретятся?
\item Из городов $A$ и $B$, расстояние между которыми $390$ км, одновременно навстречу друг другу выехали автомобили со скоростями $60$ и $70$ км/ч соответственно. На каком расстоянии от города $B$ произойдёт встреча?
\item Два автомобиля одновременно начали движение навстречу друг другу из городов, разделённых $540$ км. Первый ехал со скоростью $70$ км/ч; через $4$ часа автомобили встретились. Найдите скорость второго.
\item Два автомобиля одновременно выехали навстречу со скоростями $48$ и $62$ км/ч и встретились через $5$ часов. Найдите первоначальное расстояние между ними.
\item Расстояние между городами $550$ км. Первый автомобиль выехал со скоростью $65$ км/ч. Через $2$ часа навстречу ему выехал второй со скоростью $75$ км/ч. Через сколько часов после выезда \emph{первого} произошла встреча?
\item Расстояние между городами $720$ км. Первый автомобиль выехал со скоростью $80$ км/ч, второй --- со скоростью $60$ км/ч, причём второй отправился на $2$ часа позже. Сколько часов был в пути \emph{второй} до встречи?
\item Из городов $A$ и $B$, расстояние между которыми $620$ км, навстречу движутся автомобили. Первый выехал из $A$ со скоростью $60$ км/ч, а через $1$ час из $B$ отправился второй со скоростью $80$ км/ч. На каком расстоянии от $A$ они встретятся?
\item Между городами $690$ км. Первый автомобиль выехал из $A$ в $B$ со скоростью $75$ км/ч. Через $2$ часа ему навстречу из $B$ отправился второй. Они встретились через $4$ часа после выезда второго. Найдите скорость второго.
\item Между городами $500$ км. Первый автомобиль выехал из $A$ в $B$. Через $2$ часа из $B$ навстречу ему выехал второй автомобиль со скоростью $60$ км/ч. Они встретились в $320$ км от $A$. Найдите скорость первого автомобиля.
\item Легковой автомобиль и грузовик одновременно выехали в одном направлении, причём грузовик был впереди на $135$ км. Скорости легкового автомобиля и грузовика равны $72$ и $45$ км/ч. Через сколько часов легковой автомобиль догонит грузовик?
\end{enumerate}

\clearpage
\section*{Ответы}
\sectionrule
\begin{center}
\renewcommand{\arraystretch}{1.45}
\begin{tabular}{@{}cll@{}}\toprule
№ & Ответ & Контрольная связь\\\midrule
1 & $4$ ч & $55x+65x=480$\\
2 & $210$ км от $B$ & $\dfrac{390-x}{60}=\dfrac{x}{70}$\\[6pt]
3 & $65$ км/ч & $70\cdot4+4x=540$\\
4 & $550$ км & $48\cdot5+62\cdot5=x$\\
5 & $5$ ч & $65x+75(x-2)=550$\\
6 & $4$ ч & $80(x+2)+60x=720$\\
7 & $300$ км от $A$ & $\dfrac{x}{60}-\dfrac{620-x}{80}=1$\\[6pt]
8 & $60$ км/ч & $75\cdot6+4x=690$\\
9 & $64$ км/ч & $5x+60\cdot3=500$\\
10 & $5$ ч & $72x=135+45x$\\\bottomrule
\end{tabular}
\end{center}
\vspace{9mm}
\subidea{Итоговая памятка}
\begin{tabularx}{\linewidth}{@{}lX@{}}\toprule
Ситуация & Выбор связи\\\midrule
Едут навстречу & Пути складываются до полного расстояния: $S_1+S_2=D$.\\
Стартовали одновременно & Их время движения до встречи одинаковое.\\
Один выехал позже & Время первого больше: $t_1-t_2=\Delta t$.\\
Едут в одну сторону & Быстрый проходит путь медленного и начальный отрыв.\\
В уравнении часы и километры & Ошибка размерности: проверьте, что сравнивается.\\\bottomrule
\end{tabularx}
\vfill
\begin{center}
{\footnotesize\textcolor{softgray}{Лёвин Артём Александрович эксклюзивно для Ярослава Верещагина}}
\end{center}
\end{document}
